\documentclass[11pt,a4paper]{article}

\usepackage{fontspec}
\IfFontExistsTF{Times New Roman}{\setmainfont{Times New Roman}}{\setmainfont{Tinos}}
\IfFontExistsTF{Consolas}{\setmonofont{Consolas}[Scale=0.88]}{\setmonofont{Liberation Mono}[Scale=0.88]}
\usepackage[english]{babel}
\usepackage{geometry}
\usepackage{xcolor}
\usepackage{enumitem}
\usepackage{listings}
\usepackage{booktabs}
\usepackage{array}
\usepackage{amsmath,amssymb}
\usepackage{titlesec}
\usepackage{fancyhdr}
\usepackage{parskip}
\usepackage{tcolorbox}
\tcbuselibrary{skins}
\usepackage{tabularx}
\usepackage{lastpage}
\usepackage{hyperref}
\usepackage{multicol}

\geometry{top=1.3cm,bottom=1.3cm,left=1.65cm,right=1.65cm}

\definecolor{primaryblue}{RGB}{18,65,115}
\definecolor{secondaryteal}{RGB}{14,102,110}
\definecolor{SoftBlue}{RGB}{232,242,250}
\definecolor{CodeBg}{RGB}{248,250,252}
\definecolor{DarkGray}{RGB}{55,55,55}
\definecolor{boxbg}{RGB}{248,250,252}
\definecolor{borderblue}{RGB}{150,185,215}

\setlength{\parindent}{0pt}
\setlength{\parskip}{0pt}
\setlist[itemize]{leftmargin=1.2em,itemsep=0.3pt,topsep=0.3pt,parsep=0pt,partopsep=0pt}
\setlist[enumerate]{leftmargin=1.3em,itemsep=0.5pt,topsep=0.3pt,parsep=0pt,partopsep=0pt}
\setlist[enumerate,2]{topsep=0.1pt,itemsep=0pt,parsep=0pt,partopsep=0pt}

\lstset{
  language=Python,
  basicstyle=\ttfamily\scriptsize,
  keywordstyle=\color{blue!75!black}\bfseries,
  stringstyle=\color{red!65!black},
  commentstyle=\color{primaryblue!80!black}\itshape,
  showstringspaces=false,
  columns=fullflexible,
  keepspaces=true,
  breaklines=true,
  frame=single,
  rulecolor=\color{primaryblue!45},
  backgroundcolor=\color{CodeBg},
  xleftmargin=0.15em,
  xrightmargin=0.15em,
  aboveskip=0.10em,
  belowskip=0.10em,
  tabsize=4
}

\pagestyle{fancy}
\fancyhf{}
\setlength{\headheight}{14pt}
\fancyhead[L]{\footnotesize\textbf{DSAI1003 -- Fundamental Programming Concepts in Python}}
\fancyhead[R]{\footnotesize\textbf{Midterm Exam (Loops \& Functions, 60 Min) --- Code: 201}}
\fancyfoot[L]{\scriptsize Faculty of Data Science and Artificial Intelligence -- National Economics University}
\fancyfoot[R]{\footnotesize Page \thepage\ of 6}
\renewcommand{\headrulewidth}{0.4pt}
\renewcommand{\footrulewidth}{0.3pt}

\titleformat{\section}{\normalsize\bfseries\color{primaryblue}}{}{0pt}{}
\titleformat{\subsection}{\small\bfseries\color{primaryblue}}{}{0pt}{}
\titlespacing*{\section}{0pt}{2.0pt}{0.5pt}
\titlespacing*{\subsection}{0pt}{1.5pt}{0.4pt}

\newtcolorbox{headerbox}{
  colback=white,
  colframe=primaryblue,
  boxrule=0.7pt,
  arc=1.2mm,
  left=2.0mm,
  right=2.0mm,
  top=0.8mm,
  bottom=0.8mm
}

\newtcolorbox{answerbox}[2][]{
  enhanced,
  colback=white,
  colframe=DarkGray!45,
  boxrule=0.5pt,
  arc=1mm,
  left=2mm,
  right=2mm,
  top=0.8mm,
  bottom=0.8mm,
  title=\footnotesize\color{DarkGray}\textbf{#2},
  #1
}

\hypersetup{
  colorlinks=true,
  linkcolor=primaryblue,
  urlcolor=primaryblue,
  pdftitle={DSAI1003 - Progress Midterm Examination (Loops and Functions, 60 Minutes)},
  pdfauthor={Faculty of Data Science and Artificial Intelligence - National Economics University (NEU)}
}

\begin{document}

% =========================================================================
% PAGE 1: OFFICIAL HEADER + SCORE TABLE + ANSWER SHEET + Q1 - Q5
% =========================================================================

\noindent
\begin{minipage}[t]{0.64\textwidth}
  \raggedright\small
  \textbf{NATIONAL ECONOMICS UNIVERSITY}\\
  \textbf{COLLEGE OF TECHNOLOGY}\\
  \textbf{FACULTY OF DATA SCIENCE \& ARTIFICIAL INTELLIGENCE (FDA)}
\end{minipage}%
\hfill
\begin{minipage}[t]{0.34\textwidth}
  \raggedleft\footnotesize
  \textbf{OFFICIAL PROGRESS EXAMINATION}\\
  Academic Year: 2026--2027\\
  Semester: Fall 2026 \quad Course: \textbf{DSAI1003}
\end{minipage}

\vspace{0.08em}

\begin{center}
  {\large\bfseries\color{primaryblue} FUNDAMENTAL PROGRAMMING CONCEPTS IN PYTHON}\\[0.04em]
  {\normalsize\bfseries PROGRESS MIDTERM EXAM: LOOPS, ITERATION \& MODULAR FUNCTIONS}\\[0.04em]
  {\footnotesize\textbf{Paper Code:} \textbf{201} \quad | \quad \textbf{Duration: 60 Minutes} \quad | \quad \textbf{Difficulty:} 80\% Basic, 20\% Medium}
\end{center}

\vspace{0.04em}

\begin{headerbox}
\begin{tabular*}{\linewidth}{@{\extracolsep{\fill}}ll@{}}
  \textbf{Candidate Full Name:} \dotfill & \textbf{Student ID (MSSV):} \dotfill \\[0.10em]
  \textbf{Class / Cohort:} \dotfill & \textbf{Exam Room / Seat No.:} \dotfill
\end{tabular*}

\vspace{0.12em}
\renewcommand{\arraystretch}{1.0}
\setlength{\tabcolsep}{3.0pt}
\begin{tabularx}{\linewidth}{|l|c|c|c||c|X|}
\hline
\textbf{Section} & \textbf{\shortstack{Sec A: Quiz\\(40\%)}} & \textbf{\shortstack{Sec B: Fill-in\\(30\%)}} & \textbf{\shortstack{Sec C: Code\\(30\%)}} & \textbf{\shortstack{Total\\(100\%)}} & \centering\arraybackslash \textbf{\shortstack{Examiner\\Signature}} \\
\hline
\textbf{Max Score} & 40 pts & 30 pts & 30 pts & \textbf{100 pts} & \\[0.22em]
\cline{1-5}
\textbf{Score} & & & & & \\[0.32em]
\hline
\end{tabularx}

\vspace{0.10em}
\raggedright
\scriptsize
\textbf{Instructions:} Closed-book examination. All electronic devices and calculators are prohibited. Total score: 100 points (40\% Quiz, 30\% Fill-in, 30\% Full Code; 80\% Basic, 20\% Medium). Mark MCQs directly in the answer sheet grid. Write code solutions in designated boxes. Do not detach pages.
\end{headerbox}

\vspace{1.0mm}
\noindent
\textbf{MULTIPLE-CHOICE ANSWER SHEET} \textit{(Write letters A, B, C, or D into corresponding boxes)}:
\begin{center}
\setlength{\tabcolsep}{4.5pt}
\renewcommand{\arraystretch}{0.92}
\begin{tabular}{|c|c|c|c|c|c|c|c|c|c|c|}
\hline
\textbf{Q} & \textbf{1} & \textbf{2} & \textbf{3} & \textbf{4} & \textbf{5} & \textbf{6} & \textbf{7} & \textbf{8} & \textbf{9} & \textbf{10} \\
\hline
\textbf{Ans} & & & & & & & & & & \\[1.4mm]
\hline
\hline
\textbf{Q} & \textbf{11} & \textbf{12} & \textbf{13} & \textbf{14} & \textbf{15} & \textbf{16} & \textbf{17} & \textbf{18} & \textbf{19} & \textbf{20} \\
\hline
\textbf{Ans} & & & & & & & & & & \\[1.4mm]
\hline
\end{tabular}
\end{center}

\vspace{2.5mm}
\section*{Section A: Multiple-Choice Quiz \hfill \normalsize\normalfont(40 points --- 20 questions, 2 pts each)}
\textit{Select the single best answer. Questions marked [Basic] test core concepts; [Medium] test composite logic.}

\begin{enumerate}[label=\textbf{Q\arabic*.},leftmargin=1.8em,itemsep=0.6pt,topsep=0.2pt]

\item \textbf{[Basic]} What is the fundamental operational distinction between a \texttt{for} loop and a \texttt{while} loop in Python?
\begin{enumerate}[label=\Alph*.,leftmargin=1.8em]
  \item A \texttt{for} loop executes indefinitely until an exception occurs, whereas a \texttt{while} loop executes a fixed number of times.
  \item A \texttt{for} loop is designed for definite iteration over sequences or iterables, whereas a \texttt{while} loop performs indefinite iteration controlled by a boolean condition.
  \item A \texttt{for} loop cannot use \texttt{break} or \texttt{continue}, whereas a \texttt{while} loop supports both.
  \item A \texttt{while} loop requires an iterable object, while a \texttt{for} loop requires an explicit boolean test condition.
\end{enumerate}

\item \textbf{[Basic]} What sequence of integers is produced by evaluating \texttt{list(range(2, 14, 3))}?
\begin{multicols}{4}
\begin{enumerate}[label=\Alph*.]
  \item \texttt{[2, 5, 8, 11, 14]}
  \item \texttt{[2, 6, 10, 14]}
  \item \texttt{[2, 5, 8, 11]}
  \item \texttt{[5, 8, 11, 14]}
\end{enumerate}
\end{multicols}

\item \textbf{[Basic]} Consider the following loop intended to print integers from 1 to 5:\\
\texttt{i = 1}\\
\texttt{while i <= 5: print(i, end=" ")}\\
What happens when this program executes?
\begin{enumerate}[label=\Alph*.,leftmargin=1.8em]
  \item It displays \texttt{1 2 3 4 5} and terminates normally.
  \item It raises a \texttt{NameError} because variable \texttt{i} is not declared.
  \item It enters an infinite loop, repeatedly printing \texttt{1} because \texttt{i} is never updated inside the loop body.
  \item A \texttt{SyntaxError} is raised by Python before runtime begins.
\end{enumerate}

\item \textbf{[Basic]} In sentinel-controlled loop design, what is the specific role of a \textbf{priming read}?
\begin{enumerate}[label=\Alph*.,leftmargin=1.8em]
  \item To obtain and evaluate the initial input value before the \texttt{while} condition is tested for the first time.
  \item To compile user console input directly into executable bytecode.
  \item To prevent any runtime exceptions from ever being raised during subsequent iterations.
  \item To automatically close and flush the standard input stream.
\end{enumerate}

\item \textbf{[Basic]} What is the direct operational effect of executing a \texttt{continue} statement within a loop body?
\begin{enumerate}[label=\Alph*.,leftmargin=1.8em]
  \item It terminates the entire loop immediately and transfers control to the statement following the loop.
  \item It skips the remaining statements in the current iteration and jumps directly to the condition test/next iteration.
  \item It pauses program execution for a specified number of milliseconds.
  \item It restarts the loop from the beginning, resetting all loop variables to their initial values.
\end{enumerate}

\end{enumerate}

\newpage

% ==============================================================================
% PAGE 2: SECTION A (Q6 - Q13)
% ==============================================================================

\begin{enumerate}[label=\textbf{Q\arabic*.},leftmargin=1.8em,itemsep=0.8pt,topsep=0.3pt,resume]

\item \textbf{[Medium]} What is the final value of \texttt{total} evaluated by the loop:\\
\texttt{total, k = 0, 1}\\
\texttt{while k < 12:}\\
\texttt{\phantom{XX}if k \% 4 == 1: total += k}\\
\texttt{\phantom{XX}k += 2}
\begin{multicols}{4}
\begin{enumerate}[label=\Alph*.]
  \item \texttt{15}
  \item \texttt{25}
  \item \texttt{1}
  \item \texttt{36}
\end{enumerate}
\end{multicols}

\item \textbf{[Basic]} When finding the maximum value in an arbitrary non-empty sequence of numbers (where all values could be negative), how should the tracking variable \texttt{max\_val} be initialized?
\begin{enumerate}[label=\Alph*.,leftmargin=1.8em]
  \item Initialized to \texttt{0}, because 0 is the universal neutral integer.
  \item Initialized to \texttt{999999} to ensure all elements compare smaller.
  \item Initialized to \texttt{None} and incremented by 1 during each iteration.
  \item Initialized to the first element in the sequence, or to negative infinity (\texttt{float('-inf')}).
\end{enumerate}

\item \textbf{[Basic]} What is the output produced by the following character traversal loop?\\
\texttt{word, digits = "DataSci2026", 0}\\
\texttt{for ch in word:}\\
\texttt{\phantom{XX}if ch.isdigit(): digits += 1}\\
\texttt{print(digits)}
\begin{multicols}{4}
\begin{enumerate}[label=\Alph*.]
  \item \texttt{0}
  \item \texttt{2}
  \item \texttt{4}
  \item \texttt{11}
\end{enumerate}
\end{multicols}

\item \textbf{[Medium]} How many times is the statement \texttt{count += 1} executed in the following nested loop?\\
\texttt{count = 0}\\
\texttt{for i in range(1, 5):}\\
\texttt{\phantom{XX}for j in range(i, 5): count += 1}
\begin{multicols}{4}
\begin{enumerate}[label=\Alph*.]
  \item \texttt{10}
  \item \texttt{16}
  \item \texttt{20}
  \item \texttt{6}
\end{enumerate}
\end{multicols}

\item \textbf{[Basic]} What is the behavior of the \texttt{break} statement when executed inside an inner nested loop?
\begin{enumerate}[label=\Alph*.,leftmargin=1.8em]
  \item It terminates all enclosing loops simultaneously and halts the program.
  \item It immediately terminates only the innermost enclosing loop, leaving outer loops running.
  \item It skips only the current iteration of the inner loop without terminating it.
  \item It raises an uncatchable \texttt{StopIteration} exception.
\end{enumerate}

\item \textbf{[Basic]} What is the primary software engineering rationale for treating a Python function as a \textbf{``black box''} during system design?
\begin{enumerate}[label=\Alph*.,leftmargin=1.8em]
  \item It converts all source code lines into encrypted binary bytecode to hide algorithms from competitors.
  \item It strictly prevents functions from accepting arguments or returning results to calling programs.
  \item It mandates that all internal local variables be accessible throughout the global module namespace.
  \item The caller needs to know only the function's interface (inputs, outputs, and purpose) without needing to understand its internal implementation.
\end{enumerate}

\item \textbf{[Basic]} In Python function mechanics, what is the precise technical difference between a \textbf{parameter} and an \textbf{argument}?
\begin{enumerate}[label=\Alph*.,leftmargin=1.8em]
  \item A parameter is the variable name listed in the function definition header; an argument is the actual value or expression passed when calling the function.
  \item An argument is the variable name in \texttt{def}; a parameter is the concrete value supplied at invocation.
  \item Parameters must strictly be numeric data types, whereas arguments can only be string literals.
  \item There is no technical distinction; parameters and arguments are exact syntactic synonyms.
\end{enumerate}

\item \textbf{[Basic]} What value is returned by a Python function that completes execution without encountering an explicit \texttt{return} statement?
\begin{multicols}{4}
\begin{enumerate}[label=\Alph*.]
  \item \texttt{0}
  \item \texttt{""}
  \item \texttt{None}
  \item \texttt{False}
\end{enumerate}
\end{multicols}

\end{enumerate}

\newpage

% ==============================================================================
% PAGE 3: SECTION A (Q14 - Q20) + SECTION B INTRO & PROBLEM B1
% ==============================================================================

\begin{enumerate}[label=\textbf{Q\arabic*.},leftmargin=1.8em,itemsep=0.7pt,topsep=0.3pt,resume]

\item \textbf{[Basic]} Which of the following function definition headers will raise a \texttt{SyntaxError} in Python?
\begin{enumerate}[label=\Alph*.,leftmargin=1.8em]
  \item \texttt{def compute(val, rate=0.05, factor=1.0):}
  \item \texttt{def compute(val=10, rate=0.05, factor=1.0):}
  \item \texttt{def compute(val=10, rate, factor=1.0):}
  \item \texttt{def compute(val, rate, factor):}
\end{enumerate}

\item \textbf{[Medium]} Given \texttt{def fee(w, r=3.5, b=6.0): return w * r + b}, what is evaluated by \texttt{fee(10, b=15.0)}?
\begin{multicols}{4}
\begin{enumerate}[label=\Alph*.]
  \item \texttt{41.0}
  \item \texttt{29.0}
  \item \texttt{65.0}
  \item \texttt{50.0}
\end{enumerate}
\end{multicols}

\item \textbf{[Basic]} What is the exact \textbf{lifetime} of a local variable declared inside a Python function?
\begin{enumerate}[label=\Alph*.,leftmargin=1.8em]
  \item It is created when the function is called and deallocated immediately when the function terminates.
  \item It persists throughout the entire program execution until the Python process terminates.
  \item It is permanently stored in disk cache and reloaded across consecutive script runs.
  \item It exists only as long as its value is greater than zero.
\end{enumerate}

\item \textbf{[Basic]} What occurs when a local variable inside a function has the identical identifier as an existing global variable?
\begin{enumerate}[label=\Alph*.,leftmargin=1.8em]
  \item Python raises an immediate \texttt{AmbiguousNameError} at compile time.
  \item The local variable \textbf{shadows} (masks) the global variable within the local scope of the function.
  \item The global variable automatically overwrites the local variable upon assignment.
  \item Python merges both values into a combined two-element list.
\end{enumerate}

\item \textbf{[Medium]} What is printed by: \texttt{def inc(s): s += 25; return s} followed by \texttt{score = 60; inc(score); print(score)}?
\begin{multicols}{4}
\begin{enumerate}[label=\Alph*.]
  \item \texttt{60}
  \item \texttt{85}
  \item \texttt{25}
  \item \texttt{None}
\end{enumerate}
\end{multicols}

\item \textbf{[Basic]} When a Python function executes \texttt{return total, avg}, what data structure is actually returned to the caller?
\begin{multicols}{4}
\begin{enumerate}[label=\Alph*.]
  \item \texttt{list}
  \item \texttt{tuple}
  \item \texttt{dict}
  \item \texttt{set}
\end{enumerate}
\end{multicols}

\item \textbf{[Medium]} In professional software engineering, what is the primary purpose of \textbf{stepwise refinement}?
\begin{enumerate}[label=\Alph*.,leftmargin=1.8em]
  \item Automatically convert interpreted Python code into hardware-level microcode.
  \item Allow functions to alter global variables without declaring explicit \texttt{global} statements.
  \item Eliminate the need for loop structures by relying exclusively on infinite recursive calls.
  \item Decompose a complex problem into smaller, cohesive, independently testable functions.
\end{enumerate}

\end{enumerate}

\vspace{-1.5mm}
\section*{Section B: Coding --- Fill-in Problems \hfill \normalsize\normalfont(30 points --- 2 problems, 10 blanks)}
\textit{Write missing Python expressions directly into designated answer slots. Difficulty: 24 pts Basic (80\%), 6 pts Medium (20\%).}

\subsection*{Problem B1: Streaming Expense Tracker \& Running Total (12 points --- 4 blanks, 3 pts each [All Basic])}
Complete the program below by providing expressions for blanks \textbf{[B1.1]} to \textbf{[B1.4]}.
The utility reads individual positive expense amounts until the user enters the sentinel value \texttt{0}, computes the cumulative sum, counts items, and prints a formatted receipt.

\begin{lstlisting}
total_expense = 0.0
item_count = 0

# [B1.1] Priming read: prompt user for first expense amount and cast to float:
amount = _______________________________________________________________  # Blank [B1.1] (3 pts - Basic)

# [B1.2] While amount is greater than 0 (sentinel is 0):
while __________________________________________________________________:  # Blank [B1.2] (3 pts - Basic)
    # [B1.3] Accumulate amount into total_expense:
    total_expense = ____________________________________________________  # Blank [B1.3] (3 pts - Basic)
    item_count = item_count + 1
    amount = float(input("Enter next expense (0 to stop): "))

# [B1.4] Print final total formatted with commas and 2 decimals (e.g. 1,250.50):
print(f"Total: ${______________________} across {item_count} items")       # Blank [B1.4] (3 pts - Basic)
\end{lstlisting}

\begin{answerbox}[height=1.8cm]{Answer Slots for Problem B1}
\textbf{[B1.1]}: \dotfill \quad \textbf{[B1.2]}: \dotfill\\[1.5mm]
\textbf{[B1.3]}: \dotfill \quad \textbf{[B1.4]}: \dotfill
\end{answerbox}

\newpage

% ==============================================================================
% PAGE 4: SECTION B --- PROBLEM B2
% ==============================================================================

\subsection*{Problem B2: Modular Shopping Cart \& Discount Processor (18 points --- 6 blanks, 3 pts each)}
\textit{Difficulty: Blanks [B2.1]--[B2.4] are Basic (12 pts); Blanks [B2.5]--[B2.6] are Medium (6 pts).}

Complete the modular shopping cart processing utility below by providing expressions for blanks \textbf{[B2.1]} to \textbf{[B2.6]}.

\begin{lstlisting}
# Helper function to compute discounted price
# [B2.1] Basic: Define function header with default parameter rate=0.10:
def apply_discount(price, ______________________________________________):  # Blank [B2.1] (3 pts - Basic)
    return price * (1.0 - rate)

prices = [25.0, 40.0, 100.0, 15.0, 60.0]
total_discounted = 0.0
max_price = 0.0

# [B2.2] Basic: Traverse through each item price in the prices list:
for p in _______________________________________________________________:  # Blank [B2.2] (3 pts - Basic)
    # [B2.3] Basic: Call helper function apply_discount on p with default rate:
    disc_price = _______________________________________________________  # Blank [B2.3] (3 pts - Basic)
    total_discounted = total_discounted + disc_price
    
    # [B2.4] Basic: Track the maximum original price seen so far:
    if p > max_price:
        max_price = ____________________________________________________  # Blank [B2.4] (3 pts - Basic)

# [B2.5] Medium: Call apply_discount on total_discounted with special loyalty rate 0.05:
final_total = __________________________________________________________  # Blank [B2.5] (3 pts - Medium)

# [B2.6] Medium: Calculate average price per item safely avoiding ZeroDivisionError:
avg_price = ____________________________________________________________  # Blank [B2.6] (3 pts - Medium)

print(f"Final Total: ${final_total:.2f} | Max Item: ${max_price:.2f}")
\end{lstlisting}

\vspace{0.4mm}
\begin{answerbox}[height=1.8cm]{Answer Slots for Problem B2}
\textbf{[B2.1] (Basic):} \dotfill \quad \textbf{[B2.2] (Basic):} \dotfill\\[1.1mm]
\textbf{[B2.3] (Basic):} \dotfill \quad \textbf{[B2.4] (Basic):} \dotfill\\[1.1mm]
\textbf{[B2.5] (Medium):} \dotfill \quad \textbf{[B2.6] (Medium):} \dotfill
\end{answerbox}

\vspace{0.6mm}
% ==============================================================================
% SECTION C: IMMEDIATELY FOLLOWS SECTION B ON PAGE 4 (NO PAGEBREAK)
% ==============================================================================
\section*{Section C: Applied Full-Code Programming \& Design \hfill \normalsize\normalfont(30 points --- 3 problems)}
\textit{Difficulty: 24 pts Basic (80\%), 6 pts Medium (20\%). Problems C1 and C2 share the common athletic training scenario below with two distinct requirements. Problem C3 evaluates a sentinel loop and score aggregation.}

\begin{tcolorbox}[colback=boxbg, colframe=secondaryteal, arc=1.2mm, boxrule=0.6pt, left=2.0mm, right=2.0mm, top=0.6mm, bottom=0.6mm]
\textbf{\small Problem Scenario for C1 \& C2: Runner Fitness Activity \& Training Session Tracker}\\[0.2mm]
\footnotesize
A fitness tracking application logs a runner's performance across $N$ training sessions based on the following rules:
\begin{itemize}[itemsep=0.1pt,topsep=0.1pt]
  \item \textbf{Inputs per Session:} Duration in minutes (\texttt{duration\_min}, float) and average speed in km/h (\texttt{speed\_kmh}, float).
  \item \textbf{Formulas per Session:}
    $\text{Distance (km)} = \text{speed\_kmh} \times (\text{duration\_min} / 60.0)$; \quad
    $\text{Calories Burned (kcal)} = \text{duration\_min} \times 11.5$.
  \item Across all $N$ sessions, the application calculates total distance (km), total calories (kcal), and identifies the maximum single-session distance.
\end{itemize}
\end{tcolorbox}

\vspace{0.3mm}
\subsection*{Problem C1: Algorithmic Specification \& Pseudocode (8 points --- [All Basic])}
\textit{Requirement: Formulate structured \textbf{pseudocode} only (e.g., using \texttt{INPUT}, \texttt{COMPUTE} / \texttt{SET}, \texttt{FOR} / \texttt{WHILE}, \texttt{PRINT}). Do NOT write Python code.}

Write a complete algorithm in pseudocode that inputs the session count $N$, iterates through all $N$ sessions, computes each session's distance and calories, and determines the \textbf{total distance}, \textbf{total calories}, and the \textbf{maximum single-session distance}.

\begin{answerbox}[height=5.2cm]{Structured Pseudocode for Problem C1 (8 points --- Basic)}
\end{answerbox}

\newpage

% ==============================================================================
% PAGE 5: SECTION C --- PROBLEM C2 (MAXIMIZED CODING AREA)
% ==============================================================================

\subsection*{Problem C2: Modular Python Function Implementation (10 points --- [All Basic])}
\textit{Requirement: Write a clean, modular Python program implementing the fitness tracking engine. Define the calculation function and implement the main batch processing loop.}

Implement the program following these specifications:
\begin{enumerate}[label=\arabic*.,leftmargin=1.8em,itemsep=0.5pt]
  \item \textbf{Function \texttt{calculate\_session\_stats(duration\_min, speed\_kmh)}:}\\
    Compute and return both \texttt{distance} (float) and \texttt{calories} (float) as a tuple using the formulas from Problem Scenario C1 \& C2.
  \item \textbf{Batch Processing Loop:}
    Prompt the user for the number of sessions $N$ (int). Use a loop to iterate $N$ times. In each iteration, prompt for \texttt{duration\_min} and \texttt{speed\_kmh}, call \texttt{calculate\_session\_stats()}, and print each session's itemized distance and calories.
  \item \textbf{Accumulation \& Output Summary:}
    Tally total distance (km) and total calories (kcal). After all sessions are processed, print the summary:
\end{enumerate}

\begin{lstlisting}[basicstyle=\ttfamily\scriptsize]
Session 1: Distance = 5.00 km | Calories = 345.0 kcal
Session 2: Distance = 7.50 km | Calories = 517.5 kcal
--------------------------------------------------------
Total Sessions: 2 | Total Distance: 12.50 km | Total Calories: 862.5 kcal
\end{lstlisting}

\vspace{0.8mm}
\begin{answerbox}[height=18.0cm]{Python Implementation for Problem C2 (10 points --- Basic Modular Functions \& Loops)}
\end{answerbox}

\newpage

% ==============================================================================
% PAGE 6: SECTION C --- PROBLEM C3 (SENTINEL ITERATION & SCORE AGGREGATOR)
% ==============================================================================

\subsection*{Problem C3: Student Quiz Score Aggregator \& Performance Analyzer (12 points)}
\textit{Difficulty Breakdown: 6 points Basic + 6 points Medium. Note: Requires a sentinel-controlled loop, modular pass-check helper function, and statistical performance tracking without complex branching.}

An instructor enters student quiz scores (scale $0.0$ to $100.0$) until entering the sentinel value \texttt{-1.0}. Write a complete, concise Python program according to these specifications:

\textbf{Requirements:}
\begin{enumerate}[label=\arabic*.,leftmargin=1.8em,itemsep=0.5pt]
  \item \textbf{Modular Helper Function [3 pts --- Basic]:}\\
    Define a function \texttt{is\_passing(score, pass\_mark=50.0)} that returns \texttt{True} if $\text{score} \ge \text{pass\_mark}$, and \texttt{False} otherwise.
  \item \textbf{Sentinel Loop \& Data Aggregation [5 pts --- 3 Basic + 2 Medium]:}\\
    Use a \texttt{while} loop to prompt for student quiz scores until the sentinel \texttt{-1.0} is entered.\\
    Inside the loop, update:
    \begin{itemize}[itemsep=0.1pt]
      \item The total count of students.
      \item The running sum of all quiz scores.
      \item The count of passing students (by calling \texttt{is\_passing(score)}).
      \item The highest score recorded so far.
    \end{itemize}
  \item \textbf{Summary Report [4 pts --- Medium]:}\\
    If no student scores were entered (count is 0), print \texttt{"No scores were recorded."} and terminate.\\
    Otherwise, compute:
    $\text{average score} = \text{total sum} / \text{total count}$ \quad and \quad
    $\text{pass rate} = (\text{passing count} / \text{total count}) \times 100\%$.\\
    Display the formatted summary report (scores with 2 decimals, pass rate with 1 decimal):
\end{enumerate}

\begin{lstlisting}[basicstyle=\ttfamily\scriptsize]
========================================================
CLASS QUIZ PERFORMANCE SUMMARY
========================================================
Total Students: 5 | Passed: 4 (80.0%)
Average Score:  74.60 | Highest Score: 92.50
========================================================
\end{lstlisting}

\begin{answerbox}[height=9.8cm]{Python Implementation for Problem C3 (12 points --- 6 Basic + 6 Medium)}
\end{answerbox}

\vfill
\begin{center}
  \rule{0.5\textwidth}{0.4pt}\\[0.8mm]
  \textbf{\color{primaryblue}--- END OF EXAMINATION PAPER ---}
\end{center}

\end{document}
